\scanpage{185}
% manuscript position: top
\mspage{101}
% manuscript position: 1–4
4.3.7 (p. 77), ce qui achève la démonstration du th. 4.3,\unsure{5|6} (p. 70-71). Il reste à établir la propriété universelle énoncée au 4.3.7, que je vais rappeler
% manuscript position: left/5–13
% slanted
\margin{Devrait être incorporé en partie d) au th. 4.4.13, \unsure{ceci étant|car c'est} \unsure{un complément|conceptuellement} \unsure{important}, et pas tellement \unsure{qu'évident|évident} finalement.}

% manuscript position: 5–9
\uline{Corollaire} 4.4.16. Soit $C$ une catégorie équivalente \unsure{à une|une} petite catégorie, stable par $\varinjlim$ de cardinal $\leq\pi$, où $\pi$ est un cardinal infini donné, et considérons le foncteur d'inclusion canonique
% manuscript position: 10
\[
 C\xrightarrow{\alpha}\operatorname{Ind}_\pi(C).\tag*{(4.4.16.1)}
\]

% manuscript position: 11–16
Ce foncteur est 2-universel pour les foncteurs $f:C\longrightarrow\unsure{\mathcal E|E}$, où $\mathcal E$ est une \deleted{\unsure{$\mathcal U$-}}\inserted{$\mathcal U$-}catégorie stable par \inserted{petites} $\varinjlim$ \deleted{de cardinal $\leq\pi$}, et $f$ un foncteur commutant aux limites inductives. En d'autres termes, on a

% manuscript position: 17–20
1°) $\alpha$ possède les propriétés précédentes, i.e. $\unsure{\operatorname{Ind}_\pi(C)|\operatorname{Ind}'_\pi(C)}$ est stable par $\varinjlim$, \deleted{de cardinal $\leq\pi$} et le foncteur $\alpha$ commute \unsure{aux} $\varinjlim$ de \deleted{\unsure{de}} cardinal $\leq\pi$ ;

% manuscript position: 21–23
2°) Pour \deleted{\unsure{toute}} \inserted{$\mathcal U$-}catégorie $\mathcal E$ stable par petites limites inductives, le foncteur $g\longmapsto g\circ\alpha$
% manuscript position: 24
\[
\underline{\mathrm{Hom}}_{!}(\operatorname{Ind}_\pi C,\mathcal E)\xrightarrow{\varphi}\underline{\mathrm{Hom}}_{!\pi}(C,\mathcal E)\tag*{(4.4.16.\unsure{1|2})}
\]
% manuscript position: 25–27
est une équivalence de catégories, où l'indice $!$ (resp $!\pi$) indique la sous-catégorie \inserted{str.} pleine de la catégorie $\underline{\mathrm{Hom}}$
% manuscript position: foot
% archive sheet number: 184
